\section{Problem Definition}
\label{sec:problem_definition}

We formulate the task as exogenous-variable-aware (ExG-aware) spatio-temporal graph forecasting. This formulation aims to predict future states of a networked system by integrating historical observations, graph-structured dependencies, and external driving factors.

\subsection{Graph Representation and Spatial Dependencies}
Let $G=(V, E, A)$ denote a spatio-temporal graph, where $V=\{v_1, \ldots, v_N\}$ is the set of $N$ nodes representing physical or logical entities (e.g., wind turbines in a farm or sensors in a city). $E$ is the set of edges representing connections, and $A \in \mathbb{R}^{N \times N}$ is the adjacency matrix encoding spatial relations. In practice, $A$ can be a physical adjacency matrix, a distance-based kernel, or a learned correlation matrix. To maintain flexibility across different domains, we further define $M_G$ as a set of graph metadata, containing auxiliary information such as node coordinates, static attributes, or environmental descriptors that are not captured by the topology $A$ alone.

\subsection{Temporal Tensors and Exogenous Variables}
At a given forecasting origin $t$, we define the historical target tensor $\mathbf{X}_{t-T+1:t} \in \mathbb{R}^{T \times N \times D_x}$, where $T$ is the look-back window and $D_x$ is the dimension of the target state at each node. Our goal is to predict the future target tensor $\mathbf{Y}_{t+1:t+H} \in \mathbb{R}^{H \times N \times D_y}$ over a forecast horizon $H$.

A critical challenge in modern forecasting is the integration of exogenous variables (ExG), which we categorize into two distinct types based on their temporal availability:
\begin{itemize}
    \item \textbf{Historical ExG} ($\mathbf{E}^{h}_{t-T+1:t} \in \mathbb{R}^{T \times D_h}$): These variables represent external conditions (e.g., historical ambient temperature or humidity) that were observed alongside the target state history. They provide contextual information to help the model identify the latent state of the system during the look-back period.
    \item \textbf{Future ExG} ($\mathbf{E}^{f}_{t+1:t+H} \in \mathbb{R}^{H \times D_f}$): These represent known or forecastable future drivers (e.g., weather forecasts or scheduled operational constraints) that are available for the duration of the forecast horizon. Unlike historical variables, future ExG serves as a conditional driver for the unknown future trajectory.
\end{itemize}

\subsection{Operational Forecasting Assumptions}
To ensure the model is applicable to real-world operational settings, we enforce a strict non-leakage constraint. At inference time $t$, the model $f_{\theta}$ is only permitted to access information available up to $t$ for target states and historical ExG, and only those future ExG variables that are genuinely known or forecast by external sources (e.g., Numerical Weather Prediction models). This ensures that $\mathbf{E}^{f}$ does not contain any "future-peeking" information that would be unavailable in a deployed environment.

\subsection{General Optimization Objective}
To simplify the notation, we encapsulate all input drivers available at time $t$ into a single input tuple $\mathcal{X}_t$:
\begin{equation}
\mathcal{X}_t \triangleq \big(\mathbf{X}_{t-T+1:t}, \mathbf{E}^{h}_{t-T+1:t}, \mathbf{E}^{f}_{t+1:t+H}, A, M_G\big).
\label{eq:input_tuple}
\end{equation}
The forecasting task is defined by a parameterized model $f_{\theta}$ which maps $\mathcal{X}_t$ to the predicted horizon $\hat{\mathbf{Y}}_{t+1:t+H}$. For domains where physical laws or operational limits must be satisfied, the model can be decomposed into an unconstrained spatio-temporal backbone $\mathcal{H}_{\theta}$ and a constraint operator $\mathcal{C}_{\Omega}$ parameterized by domain knowledge $\Omega$:
\begin{equation}
\hat{\mathbf{Y}}_{t+1:t+H} = \mathcal{C}_{\Omega}\big(\mathcal{H}_{\theta}(\mathcal{X}_t)\big).
\label{eq:constrained_forecast}
\end{equation}
The operator $\mathcal{C}_{\Omega}$ ensures that the final output adheres to feasibility requirements such as non-negativity or capacity limits.

Given a training dataset $\mathcal{D}$, the model parameters $\theta$ are optimized by minimizing the Mean Absolute Error (MAE) between the predictions and the ground truth:
\begin{equation}
\mathcal{L}(\theta) = \frac{1}{|\mathcal{D}| \cdot H \cdot N} \sum_{(\mathcal{X}, \mathbf{Y}) \in \mathcal{D}} \sum_{h=1}^{H} \sum_{n=1}^{N} | f_{\theta}(\mathcal{X})_{h,n} - \mathbf{Y}_{h,n} |,
\label{eq:optimization_objective}
\end{equation}
where the $L_1$ loss provides robustness against outliers in the spatio-temporal data. This objective provides a general training framework that can be easily extended with additional task-specific regularizers in later stages of the methodology.
